% Borrador de sección para el informe (LaTeX, prosa fluida).
% Documenta la exploración "¿puede la no-linealidad superar los baselines aditivos?".
% Todos los números son de las corridas de validación honesta (regla de oro + bootstrap
% pareado cluster-por-línea). Reemplazar los \pearson por los links de wandb al integrar.

\section{¿Aporta la no-linealidad sobre los baselines aditivos?}
\label{sec:no-linealidad}

Establecidos los baselines por ambiente —GBLUP para el ambiente óptimo y el LASSO
baggeado para sequía—, la pregunta natural es si una red neuronal, u otro modelo capaz
de capturar estructura no-aditiva, puede superarlos explotando interacciones entre
marcadores. La hipótesis es atractiva: si existiera epistasis o dominancia relevante para
el rendimiento, un modelo aditivo la dejaría sobre la mesa. Esta sección la pone a prueba
de forma sistemática y concluye que, para estos datos, la no-linealidad no ofrece una
mejora robusta: la arquitectura genética del rendimiento es aditiva en ambos ambientes y
los baselines ya operan en el techo de predictibilidad.

\subsection{Protocolo de validación}

Toda comparación se juzga con dos folds temporales y un criterio deliberadamente estricto
que llamamos la \emph{regla de oro}. El primer fold es el análogo de despliegue: se
entrena con 2017+2018 y se valida en 2019, replicando la situación de predecir el año
siguiente. El segundo es un fold de robustez tipo \emph{leave-one-year-out} que entrena
con 2017+2019 y valida en 2018. Un método se considera superador solo si mejora al
baseline en \emph{ambos} folds; el año 2020 permanece sellado. Las diferencias de
correlación de Pearson se reportan con un bootstrap pareado que remuestrea clusters por
línea, de modo que el intervalo de confianza refleja la incertidumbre real a nivel de
genotipo y no de observación. Este diseño es intencionalmente severo: en régimen
$p \gg n$ es fácil obtener ganancias espurias que solo aparecen en un fold —el clásico
\emph{winner's curse} discutido en \citep{murphy2022}— y la regla de oro las descarta.

\subsection{Métodos evaluados}

La exploración cubrió las principales familias con las que, en teoría, se podría capturar
señal no-aditiva u ortogonal. Entre las redes y kernels se probaron perceptrones
multicapa sobre los marcadores crudos, redes residuales que parten de la predicción del
ganador, kernels gaussianos (RKHS, la versión no-lineal de GBLUP), esquemas multi-kernel
que combinan un kernel aditivo con uno gaussiano, máquinas de factorización para epistasis
de pares, autoencoders de denoising con salida sigmoide y entropía cruzada binaria,
aumentación de datos con un C-$\beta$-VAE, y redes multitarea y multi-ambiente. Entre los
enfoques de combinación se evaluó el \emph{stacking} de las predicciones de GBLUP y LASSO
—tanto in-sample como con predicciones \emph{out-of-fold}—, la mezcla de pesos fijos, la
transferencia entre ambientes (usar la predicción genómica del rendimiento en óptimo como
covariable para sequía), el uso de los caracteres secundarios altura de planta y días a
floración como predictores, y árboles con \emph{boosting} en sus variantes directa y
residual. Finalmente se barrió el \emph{elastic net} para verificar si una penalización
mixta $L_1$+$L_2$ captura simultáneamente una componente densa y una \emph{sparse}.

\subsection{Resultado: ningún método supera de forma robusta}

Ninguna de las familias anteriores mejora a su baseline en los dos folds. El patrón es
sistemático y esclarecedor. En el ambiente óptimo, cualquier modelo que introduzca
\emph{sparsity} —LASSO, elastic net, árboles— derrumba la predicción en el fold de
despliegue (por ejemplo, el elastic net cae de $0{,}31$ a valores de $0{,}13$--$0{,}17$ en
2019), aunque supere a GBLUP en el fold de 2018. Esta inversión por fold no es explotable:
2019 exige la densidad de GBLUP y 2018 premia la \emph{sparsity}, de modo que ninguna
combinación fija ni ningún modelo mixto gana en ambos. La prueba más nítida proviene del
esquema multi-kernel, cuya validación interna asigna peso nulo al kernel gaussiano y
recupera exactamente GBLUP: no hay varianza no-aditiva que ponderar. En sequía, donde la
señal es \emph{sparse}, el LASSO baggeado resulta ser el techo; el elastic net empata a lo
largo de todo el rango de $l_1$ (las diferencias caen dentro del ruido del bootstrap), lo
que indica que no existe una cola poligénica densa que una componente $L_2$ pueda
aprovechar.

El \emph{stacking} merece mención aparte por lo instructivo de su fracaso. Combinar las
predicciones de GBLUP y LASSO con un meta-modelo entrenado sobre las predicciones
in-sample colapsa, porque el combinador sobre-confía en el modelo que mejor sobreajusta el
train y en validación pondera mal al que realmente generaliza. Corregido con predicciones
\emph{out-of-fold}, el \emph{stacking} honesto (mínimos cuadrados no negativos o Ridge)
pierde en todos los folds, y el escaneo de pesos fijos ubica su óptimo en el extremo —la
base pura— confirmando que las predicciones densa y \emph{sparse} no aportan información
complementaria una sobre la otra.

\subsection{El único indicio positivo y su refutación}

Un solo experimento produjo una mejora aparente y merece ser documentado porque ilustra la
importancia del protocolo. La transferencia entre ambientes —predecir genómicamente el
rendimiento en condiciones óptimas de cada línea, que por su mayor heredabilidad captura el
potencial genético general, y usarlo como covariable para sequía— dio, con una red
meta-aprendiz, una ganancia de $+0{,}017$ en Pearson sobre el LASSO en el fold de
despliegue de 2019 (intervalo $[+0{,}010,\,+0{,}025]$, estable sobre dieciséis semillas).
Fue el único caso en que un método mejoró un baseline sin dañar el otro fold. Sin embargo,
la red se entrenaba sobre las predicciones \emph{in-sample} de los modelos base. Al repetir
exactamente el mismo modelo pero entrenándolo sobre predicciones \emph{out-of-fold} —el
único procedimiento honesto y efectivamente desplegable— la ganancia no solo desaparece
sino que se vuelve negativa ($-0{,}039$ en 2019 y $-0{,}29$ en el fold de 2018). La mejora
era, por lo tanto, optimismo in-sample de las predicciones base, no señal transferible. El
episodio subraya por qué la selección honesta y la validación en más de un fold son
indispensables.

\subsection{Interpretación}

La convergencia de todos estos resultados admite una lectura única, y existe una razón
biológica que la explica de raíz. Los marcadores están codificados de forma binaria, lo que
indica que las líneas son endocriadas y homocigotas: no hay heterocigotas en el panel. En
consecuencia, la varianza de dominancia es estructuralmente nula —la dominancia requiere
genotipos heterocigotas— y solo pueden existir componentes aditiva y epistática. Como la
epistasis se descartó empíricamente por varias vías (máquinas de factorización, kernel
gaussiano con peso nulo, árboles), queda esencialmente la componente aditiva. Que GBLUP y el
LASSO —ambos modelos aditivos— constituyan el techo no es entonces un accidente de estos
datos sino una consecuencia esperable de trabajar con líneas endocriadas. Un fenómeno
secundario refuerza el punto: en el ambiente óptimo el rendimiento relativo de los métodos
se invierte entre folds —los métodos locales o \emph{sparse} superan a GBLUP cuando se
valida en 2018 pero se derrumban en 2019—, y esa inversión no es explotable porque la señal
que indicaría en qué régimen se está (el grado de parentesco de las líneas a validar con el
conjunto de entrenamiento) no se transfiere de forma fiable: un predictor adaptativo que
pondera GBLUP contra vecinos según el parentesco, con pesos aprendidos honestamente fuera de
muestra, aprende la regla equivocada y pierde en el fold de despliegue.

El rendimiento de grano es, en este conjunto de datos y para ambos ambientes, un carácter de
arquitectura aditiva.
En el ambiente óptimo domina un régimen infinitesimal denso para el cual GBLUP, que es
\emph{ridge regression} sobre la matriz de parentesco genómico, es cercano al óptimo de
Bayes bajo un \emph{prior} gaussiano \citep{bishop2024}. En sequía la señal se concentra en
pocos marcadores y el LASSO baggeado, estimador \emph{maximum a posteriori} de un
\emph{prior} de Laplace, la modela mejor, con el \emph{bagging} reduciendo su varianza. En
ninguno de los dos casos el residuo del ganador contiene estructura no-aditiva que una red
pueda recuperar: cuando se la fuerza a modelarlo, agrega varianza y empeora la validación.
El resultado no es un fracaso de implementación —se probaron del orden de quince familias
de métodos con búsqueda de hiperparámetros y múltiples semillas— sino una propiedad de la
señal. La mejora sustantiva del proyecto, en consecuencia, no proviene de la no-linealidad
sino del cambio de \emph{prior} apropiado al ambiente: el LASSO supera a GBLUP en sequía en
torno a un $63\%$ en Pearson, y ese salto, aditivo y \emph{sparse}, es el hallazgo central.
